\section{Conclusion}
AC--MOT is best understood as an optimization-driven detector operating-point
interface rather than as a single timeless heuristic. V1 was the first formal
quality-oriented optimization and produced a quality profile with uncertainty.
V2 then made identity and deployment cost explicit through a feasible Pareto
search, showing why a lower-IDS/faster profile need not be the strongest HOTA
profile.
Historical UAVDT transfer, U2MOT on VisDrone, and SparseTrack on MOT17 further
showed that profile choice and tracker host matter, including negative and
near-neutral cases.

The modern frozen AC--MOT v3 phase tightened the protocol around a stronger
detector, a crowd-selected SCI with causal guards, a finite action map, two
tracker hosts, paired bootstrap evidence, T4 instrumentation, and no-retuning
UAVDT transfer. The modern results support positive changes over the declared
\texttt{r1536\_n70} baseline on ByteTrack and OATrack, with stronger ByteTrack
effects and mixed OATrack transfer.

The final conclusion is that AC--MOT supports objective-driven detector
operating-point control with tracker- and dataset-dependent outcomes. The
current record does not establish universal tracker independence,
timing-specific causality, or broader detector generalization; these remain
questions for explicitly designed future experiments.
